%! TeX root = main.tex
% ============================================================
% CONCLUSION GÉNÉRALE ET PERSPECTIVES
% ============================================================

\chapter*{Conclusion Générale et Perspectives}
\addcontentsline{toc}{chapter}{Conclusion Générale et Perspectives}

Ce travail a permis de concevoir et de développer une application web
dédiée à la gestion des travaux dirigés destinés aux candidats aux
examens du CEP, du BEPC et du Baccalauréat de la commune de Cotonou.
Les fonctionnalités mises en œuvre --- gestion du cycle de vie des TDs,
suivi des statuts, traitement des paiements, génération de rapports et
tableau de bord en temps réel --- apportent une réponse numérique
concrète aux insuffisances des méthodes manuelles encore en vigueur
dans ces services administratifs. Le projet met en lumière le potentiel
de la transformation numérique dans la gestion publique locale, en
particulier dans le contexte béninois où les solutions digitales
adaptées aux réalités administratives locales restent insuffisantes.

La structuration du système autour de trois profils distincts
--- administrateur, enseignant et comptable --- garantit une
séparation claire des responsabilités, une traçabilité complète
des opérations et une sécurisation rigoureuse des données.
Certaines limites subsistent néanmoins, notamment la dépendance
à une connexion internet stable, l'absence d'une application mobile
native et le recours à un serveur SMTP externe pour les notifications
par e-mail. Ces limites ne remettent pas en cause la pertinence
de la solution mais constituent des pistes d'amélioration
concrètes pour les versions futures.

\bigskip
\noindent\textbf{Les perspectives d'évolution de l'application
s'articulent autour des axes suivants :}

\begin{itemize}
    \item intégrer un module de paiement mobile directement dans
    la plateforme, en s'appuyant sur des solutions adaptées
    au contexte béninois telles que MTN Mobile Money ou Moov Money,
    afin d'automatiser entièrement le processus de règlement
    des enseignants ;

    \item développer une application mobile native (Android et iOS)
    afin d'améliorer l'accessibilité de la plateforme et de
    proposer des notifications push en temps réel à chaque acteur
    du système ;

    \item enrichir le module statistique du tableau de bord
    avec des graphiques d'évolution, des comparaisons inter-périodes
    et des exports automatisés en PDF, pour renforcer les capacités
    d'analyse et de reporting des décideurs municipaux ;

    \item concevoir une architecture multi-tenant permettant
    d'étendre la solution à d'autres communes béninoises, chacune
    disposant de son espace isolé et sécurisé au sein d'une
    infrastructure partagée.
\end{itemize}

En définitive, cette application démontre que l'innovation numérique
appliquée à l'administration publique locale peut générer une réelle
valeur ajoutée, améliorer la transparence des opérations financières
et ouvrir la voie à une modernisation durable des services municipaux
au Bénin.